\subsection{Key Findings and Implications for Methodology}
\label{sec:key_findings}
The dataset analysis establishes a controlled starting point for the five-model
comparison. Its main findings translate into the following experimental
decisions.

\begin{enumerate}
\item \textbf{Class balance supports an unweighted baseline.}
Each class contributes 5,400 training, 600 validation, and 1,000 test images.
No resampling or class weighting is introduced to address frequency imbalance.
Equal counts do not guarantee equal recognition difficulty, so evaluation
must include per-class precision, recall, and F1 alongside accuracy and macro-F1.

\item \textbf{Background prevalence must be separated from image validity.}
Exactly \ZeroPercent\% of training pixels are zero, but no training image
is entirely blank. Brightness and coverage distributions overlap across
classes. The baseline therefore keeps the grayscale intensities and every
image rather than applying an unsupported brightness threshold or background
removal rule.

\item \textbf{Spatial structure motivates different model representations.}
Class means retain broad silhouettes, while variability maps and individual
examples show changes that averaging conceals. Vectors provide the Linear
and MLP baselines; the CNN retains the image grid; the LSTM reads ordered
rows; and the Transformer receives ordered patches. All layouts retain the
same 784 standardized values. Architecture comparisons must still account
for differences in parameter count and computation before attributing a
performance gap to representation alone.

\item \textbf{Prototype similarity gives a specific error-analysis question.}
Coat and Shirt have the highest off-diagonal prototype cosine similarity,
approximately \MostSimilarCosine{}. Later confusion matrices should check
whether these and other upper-body categories are frequent directional
errors, and example galleries should examine the corresponding mistakes.
Similarity of averages alone does not predict which model will perform best
or establish that any supplied label is incorrect.

\item \textbf{One fitted preprocessing pipeline supports a consistent comparison.}
The seed-42 split, cached indices, training-only normalization, and verified
input ordering are shared by all architectures. Validation loss selects
weights, and the selected checkpoint is evaluated on the official test set.
Any change to augmentation or representation must be recorded as a separate
experimental condition rather than silently changing the baseline.
\end{enumerate}

The analysis has limits. Exact-image hashing cannot exclude near duplicates
or related products, and a few sampled images cannot establish the frequency
of an appearance. Class averages hide individual variation. The present
evidence also does not measure training variability or generalization to
high-resolution, colored, or unconstrained product photographs. Those limits
remain relevant when interpreting later results.

Section~3 can now define the five architectures, their mathematical
operations, loss, training strategy, and evaluation metrics using the input
shapes established in Section~2.4. Section~4 can describe how the shared
dataset views, batching, random-state control, device transfer, and checkpoint
handling implement that protocol. Training histories, predictive results,
timing measurements, and error analysis belong to the later experimental
sections; the EDA provides hypotheses and preprocessing decisions for those
experiments rather than substituting for them.
